\documentclass{article}

% Recommended packages for figures and better typesetting
\usepackage{microtype}
\usepackage{graphicx}
\usepackage{subcaption}
\usepackage{booktabs}

% hyperref makes hyperlinks in the resulting PDF
\usepackage{hyperref}

% Attempt to make hyperref and algorithmic work together better
\newcommand{\theHalgorithm}{\arabic{algorithm}}

% Use the following line for the initial blind version submitted for review
\usepackage{icml2026}

% For preprint, use:
% \usepackage[preprint]{icml2026}

% If accepted, instead use the following line for the camera-ready submission:
% \usepackage[accepted]{icml2026}

\usepackage{amsmath}
\usepackage{amssymb}
\usepackage{mathtools}
\usepackage{amsthm}

% if you use cleveref
\usepackage[capitalize,noabbrev]{cleveref}

\icmltitlerunning{GPU Architecture Evolution: A Bottleneck-Driven Survey}

\begin{document}

\twocolumn[
\icmltitle{GPU Architecture Evolution Under ML Workload Pressure:\\
A Causal, Bottleneck-Driven Analysis from Fermi to Rubin}

\icmlsetsymbol{equal}{*}

\begin{icmlauthorlist}
\icmlauthor{Heng Xiong}{equal,fudan}
\icmlauthor{Claude}{equal,anthropic}
\end{icmlauthorlist}

\icmlaffiliation{fudan}{Fudan University, Shanghai, China}
\icmlaffiliation{anthropic}{Anthropic}

\icmlcorrespondingauthor{Heng Xiong}{25300130028@m.fudan.edu.cn}

\icmlkeywords{GPU Architecture, Deep Learning, Hardware Acceleration, Performance Bottlenecks}

\vskip 0.3in

\begin{abstract}
We present a systematic analysis of GPU architecture evolution from NVIDIA Fermi (2010) through Rubin (2026), examining how each architectural transition directly addresses performance bottlenecks created by emerging machine learning workloads. We identify a clear causal chain: Fermi addressed HPC memory bandwidth, Volta introduced Tensor Cores for matrix multiplication throughput, Ampere added structured sparsity and TF32, Hopper optimized for transformer attention with FP8, and Blackwell targets inference efficiency with FP4. This survey provides both a historical perspective and predictive framework for understanding future GPU evolution as ML workloads continue to scale.
\end{abstract}
]

\printAffiliationsAndNotice{}

\section{Introduction}

The explosive growth of machine learning, particularly deep learning, has fundamentally transformed GPU architecture over the past 15 years. NVIDIA's CUDA platform (2007) initiated the GPGPU computing revolution, positioning GPUs as general-purpose parallel processors---but the true architectural revolution began with three converging forces: the 2012 AlexNet breakthrough demonstrating deep learning's potential \cite{krizhevsky2012imagenet}, the scaling hypothesis showing larger models consistently improve \cite{kaplan2020scaling}, and the transformer's emergence as the dominant paradigm \cite{vaswani2017attention}. Before these workloads, GPUs served HPC (fluid dynamics, molecular dynamics) and graphics; early GPGPU architectures lacked ECC memory, true cache hierarchies, and double-precision support required for scientific computing.

This survey examines GPU architecture evolution through a \textit{bottleneck-driven lens}---each generation's design directly addresses the dominant performance bottleneck of its era. We analyze seven GPU generations:

\begin{itemize}
\item \textbf{Fermi (2010)}: Pre-deep learning HPC optimization
\item \textbf{Volta (2017)}: First Tensor Cores for matrix multiplication
\item \textbf{Turing (2018)}: Inference specialization with INT8
\item \textbf{Ampere (2020)}: Sparsity and TF32 for training
\item \textbf{Hopper (2022)}: Transformer-specific FP8 acceleration
\item \textbf{Blackwell (2024)}: Inference scaling with FP4
\item \textbf{Rubin (2026)}: Energy efficiency and extreme-scale inference
\end{itemize}

For each generation, we examine ten fundamental questions spanning workload drivers, computational primitives, memory hierarchies, numerical precision, and performance scaling.

\section{Fermi (2010): Establishing the Foundation}

\subsection{Workload Context and Bottleneck}

Fermi (GF100/GF104) emerged as scientific computing demanded GPU acceleration, but existing architectures lacked reliability and precision \cite{nvidia2009fermi}. The dominant bottleneck was \textit{memory bandwidth for irregular HPC codes}---previous GPUs (G80, GT200) delivered high peak FLOPS but suffered on applications with poor spatial locality. Deep learning remained nascent; AlexNet would not arrive until 2012.

\subsection{Computational Model and Key Innovations}

Fermi's dominant unit was the \textbf{CUDA Core}---a scalar FP32/FP64 ALU with dual-issue capability. Each of 16 SMs contained 32 CUDA Cores (512 total), with FP64 running at 1/2 FP32 rate. This was a \textit{fully general-purpose} architecture with no specialized ML units.

\textbf{Innovations addressing the HPC bottleneck}:
\begin{itemize}
\item \textbf{True cache hierarchy}: Configurable L1 (16/48KB) + 768KB unified L2; revolutionary over prior all-programmer-managed shared memory
\item \textbf{ECC + IEEE 754-2008}: Full ECC on GDDR5, correct rounding, denormals, NaN handling
\item \textbf{Dual warp schedulers}: Two-way ILP per SM; 48 warps (1536 threads) with GigaThread hardware context switching
\end{itemize}

\textbf{Memory and bottleneck}: 384-bit GDDR5 at 177 GB/s. Despite the cache hierarchy, memory bandwidth remained the limiting factor---at 177 GB/s and 1.03 TFLOPS FP32 (peak), the compute-to-bandwidth ratio was only $\sim$5.8 FLOPs/byte, insufficient for streaming HPC.

\textbf{Legacy}: Fermi resolved the HPC bottleneck and established invariants (SM structure, warp-based SIMT, shared memory programming) that persist across all subsequent generations. It exposed the \textit{next} constraint: matrix multiplication throughput for the emerging deep learning revolution.

\section{Volta (2017): The Tensor Core Revolution}

\subsection{Workload Context and Bottleneck}

Between Fermi and Volta, deep learning exploded: AlexNet (2012), ResNets (2015), and early attention mechanisms previewed transformers \cite{bahdanau2014neural}. Training times stretched to weeks on Pascal P100, bottlenecked by \textit{matrix multiplication throughput}. Volta (GV100) addressed this directly with \textbf{Tensor Cores}---specialized MMA units delivering 12x speedup over Pascal for mixed-precision training \cite{nvidia2017volta}.

\subsection{Tensor Cores, Memory, and Precision}

\begin{table}[h]
\centering
\footnotesize
\begin{tabular}{ll}
\toprule
Attribute & Specification \\
\midrule
Operation & $D = A \times B + C$ (4$\times$4$\times$4 MMA) \\
Precision & FP16 input, FP32 accumulate \\
Tensor Cores/SM & 8 (640 total, 80 SMs) \\
Peak throughput & 125 TFLOPS (8x CUDA: 15.7 TF) \\
\bottomrule
\end{tabular}
\caption{Volta V100 Tensor Core specifications showing 8x advantage over CUDA Cores.}
\end{table}

This marks an \textit{architectural phase transition}: general-purpose scalar ALUs give way to specialized matrix units. The V100 retained 5120 CUDA Cores for non-matrix work, but Tensor Cores dominated ML performance.

\textbf{Memory}: HBM2 with 4096-bit interface, 900 GB/s (5x Fermi), 16/32GB. The compute-to-bandwidth ratio reached $\sim$139 FLOPs/byte---workloads became compute-bound rather than bandwidth-bound. Independent thread scheduling enabled fine-grained synchronization.

\textbf{Precision}: Volta established \textbf{mixed-precision training} as the standard \cite{micikevicius2017mixed}---FP16 inputs, FP32 accumulation, loss scaling to prevent gradient underflow---while retaining FP64 (7.8 TFLOPS) for HPC.

\textbf{Legacy}: Tensor Cores, mixed precision, and HBM become permanent fixtures. \textbf{Next bottleneck}: attention's $O(n^2)$ complexity and the need for lower inference precision.

\section{Turing (2018): Inference and Ray Tracing}

\subsection{Workload Context and Divergent Optimization}

Turing (TU102/TU104) represents a \textit{fork in GPU evolution}: Volta targeted data center training; Turing optimized for \textbf{inference deployment} and \textbf{real-time ray tracing} \cite{nvidia2018turing}. Training is throughput-oriented (large batches); inference is latency-critical (small batches, real-time constraints).

\subsection{Dual Specialization and Memory}

Turing's 70\% specialized architecture introduced two new unit types:
\begin{itemize}
\item \textbf{2nd-gen Tensor Cores}: INT8 (65 TOPS, 8x FP16 throughput), INT4 (130 TOPS, 16x), and FP16 (8.1 TFLOPS)---enabling post-training quantization with $<$1\% accuracy loss via QAT
\item \textbf{RT Cores}: Ray-triangle intersection and BVH traversal for real-time ray tracing (DLSS)
\end{itemize}

\textbf{Memory}: GDDR6 (616 GB/s, 11GB) rather than HBM2, reflecting consumer GPU economics. Lower bandwidth makes inference latency the bottleneck: large batches become memory-bound; small batches underutilize Tensor Cores.

Turing excels at edge AI (YOLO, DLSS super-resolution) but is not optimal for large model training. \textbf{Next bottleneck}: training throughput for GPT-scale models, sparsity exploitation, and cloud multi-tenancy efficiency.

\section{Ampere (2020): The Sparsity and TF32 Era}

\subsection{Workload Context and Bottleneck}

Ampere (GA100/A100) arrived as GPT-3 (175B parameters, 2020) \cite{brown2020language} created massive compute demands. Three simultaneous bottlenecks: (1) training throughput needing 10--100x more compute, (2) pruned networks lacking hardware sparsity support, and (3) single-tenant cloud GPUs with $\sim$30\% utilization.

\subsection{Precision Diversity, Sparsity, and MIG}

A100's 3rd-gen Tensor Cores introduced \textbf{precision diversity} \cite{nvidia2020ampere}:
\begin{itemize}
\item \textbf{TF32}: FP32 range + FP16 precision (10-bit mantissa)---automatic 8x speedup for legacy FP32 code with no code changes
\item \textbf{BF16}: Wider dynamic range than FP16, better numerical stability for large transformers
\item \textbf{2:4 structured sparsity}: Of every 4 values, exactly 2 must be zero; Tensor Core compresses and doubles effective throughput (312 $\to$ 624 TFLOPS sparse) with $<$1\% accuracy loss \cite{mishra2021accelerating}
\end{itemize}

\textbf{Multi-Instance GPU (MIG)}: Partitions one GPU into up to 7 isolated instances with dedicated SMs, memory controllers, and L2 partitions---improves cloud utilization from $\sim$30\% to $\sim$90\%.

\textbf{Memory}: HBM2e at 40GB (1.6 TB/s) or 80GB (2.039 TB/s); 40MB L2 (6.7x V100). GPT-3 175B fits on 8$\times$A100 with model parallelism. Async copy and hardware task graphs overlap communication with computation.

\textbf{Performance}: TF32 at 156 TFLOPS; FP16 at 312 TFLOPS (2.5x over V100's 125 TFLOPS FP16); NVLink 3.0 at 600 GB/s (2x V100). \textbf{Remaining bottleneck}: Transformer attention's $O(n^2)$ complexity---at 2048 tokens, attention consumes 40\% of training time.

\section{Hopper (2022): The Transformer Engine}

\subsection{Workload Context and Bottleneck}

By 2022, models reached trillion-parameter scale: PaLM 540B and Megatron-Turing NLG 530B trained on thousands of GPUs for months \cite{chowdhery2022palm}. Three bottlenecks: (1) transformer attention consuming 40--60\% of training time, (2) FP16 hitting numerical instability at trillion-parameter scale, and (3) all-reduce communication limiting multi-GPU scaling.

\subsection{Transformer Engine, FP8, and Architecture}

Hopper (GH100/H100) introduced the \textbf{Transformer Engine}---hardware + software co-design for attention \cite{nvidia2022hopper}:
\begin{itemize}
\item \textbf{FP8 formats}: E4M3 (4-bit exp, 3-bit mantissa, high precision) and E5M2 (5-bit exp, 2-bit mantissa, wide range)
\item \textbf{Automatic per-layer switching}: Tensor statistics guide FP8/FP16 selection; $<$0.5\% accuracy degradation; 2x throughput over FP16, delivering 2000 TFLOPS with FP8 versus 1000 TFLOPS with FP16
\end{itemize}

\textbf{Thread block clusters}: SM-to-SM distributed shared memory---7x faster than global memory---enables large-matrix tiling across SMs for massive GEMM.

\textbf{Memory}: HBM3 at 80GB, 3.35 TB/s (2x A100); 50MB L2; Tensor Memory Accelerator (TMA) frees threads from addressing overhead. Compute-to-bandwidth ratio reaches $\sim$597 FLOPs/byte (FP8). DPX instructions deliver 7x speedup for genomics (Smith-Waterman) and robotics routing.

\textbf{Communication}: NVLink 4.0 at 900 GB/s per GPU; 256-GPU NVLink Switch at 57.6 TB/s all-to-all.

\textbf{Performance}: FP8 Tensor at 2000 TFLOPS (4000 sparse)---6.4x over A100 FP16; 9x training speedup vs.\ A100. \textbf{Remaining bottleneck}: Attention's $O(n^2)$ persists; FlashAttention helps algorithmically \cite{dao2022flashattention}.

\section{Blackwell (2024): Inference at Scale}

\subsection{Workload Context and Bottleneck}

At GPT-4 scale, a new bottleneck emerged: \textbf{inference deployment at scale}. Serving trillion-parameter models to millions of users demands extreme throughput (tokens/second), low latency ($<$100ms to first token), and cost efficiency (\$/token)---constraints fundamentally different from training.

\subsection{Multi-Chip Architecture, FP4, and Scale}

Blackwell (GB100/GB200) introduces \textbf{dual-die design}: two reticle-limited dies at 10 TB/s chip-to-chip, yielding 208 billion transistors (2.6x H100) \cite{nvidia2024blackwell}.

\textbf{Transformer Engine v2 with FP4}:
\begin{itemize}
\item NVFP4 with micro-tensor scaling: 2x memory capacity and 2x throughput vs.\ FP8
\item Enables 400B parameter models in budgets previously limited to 200B
\end{itemize}

\begin{table}[h]
\centering
\footnotesize
\begin{tabular}{ll}
\toprule
Metric & Value \\
\midrule
FP8/FP4 (dense) & 5/10 PFLOPS \\
FP8/FP4 (sparse) & 10/20 PFLOPS \\
Memory (HBM3e) & 192 GB, 8 TB/s \\
NVL72 bandwidth & 130 TB/s \\
\bottomrule
\end{tabular}
\caption{Blackwell (GB200) key specifications (208B transistors, dual-die).}
\end{table}

\textbf{Memory}: HBM3e at 192GB per GPU (2.4x H100), 8 TB/s; hardware LZ4/Snappy/Deflate decompression; 900 GB/s Grace-Blackwell CPU-GPU coherent link enables 100K+ token inference without CPU spill.

\textbf{NVL72}: 72 GPUs in a single NVLink domain, 130 TB/s aggregate---the rack becomes the unit of computation. RAS Engine applies AI-powered fault prediction for data center reliability.

\textbf{Performance}: NVL72 delivers 30x inference speedup and 25x lower cost vs.\ H100 for GPT-4 scale models. \textbf{Remaining bottleneck}: Test-time compute scaling for reasoning models.

\section{Rubin (2026): Energy Efficiency at Extreme Scale}

\subsection{Workload Context and Bottleneck}

By 2026, multi-trillion parameter models (10T+), real-time multimodal AI, and extended reasoning push compute demands to new extremes. The dominant bottleneck shifts to \textbf{energy efficiency and inference token cost}: data centers approach physical megawatt-scale power limits, making performance-per-watt and cost-per-token the critical metrics for sustainable AI scaling.

\subsection{Architectural Response}

Rubin, officially announced at CES 2026, is NVIDIA's first architecture explicitly designed for energy-constrained extreme-scale AI \cite{nvidia2024rubin}. Built on TSMC 3nm process with 336 billion transistors (dual reticle design), Rubin delivers:

\textbf{6th-gen Tensor Cores with NVFP4}: 50 PFLOPS NVFP4 inference performance (5x Blackwell); 35 PFLOPS NVFP4 training (3.5x Blackwell); extended FP4 precision support with micro-tensor scaling; enhanced FP8 capabilities; hardware-accelerated adaptive compression; improved confidential computing for secure multi-tenant AI.

\textbf{HBM4 memory}: 288 GB capacity per GPU (1.5x Blackwell); 22 TB/s bandwidth (2.75x Blackwell's 8 TB/s); compute-to-bandwidth ratio of $\sim$2273 FLOPs/byte at NVFP4 precision.

\textbf{NVLink 6}: 3.6 TB/s bidirectional bandwidth per GPU (2x NVLink 5); electrical interconnect (not optical); NVL72 configuration with 260 TB/s total scale-up bandwidth across 72 GPUs.

\textbf{Rubin Platform}: Features the next-generation \textbf{Vera} Arm-based CPU for orchestration; next-generation SuperNICs for cluster networking; advanced DPUs for data processing; enhanced RAS Engine for data center reliability.

Rubin achieves 10x lower inference token cost versus Blackwell, directly addressing the energy efficiency and deployment cost bottlenecks at extreme AI scale. Full production began Q1 2026, with products shipping H2 2026. \textbf{Remaining bottleneck}: Beyond 3nm, Moore's Law scaling slows---photonic interconnects or processing-in-memory may be required for continued progress.

\section{Cross-Generational Analysis}

\subsection{Bottleneck Evolution Chain}

\begin{equation}
\begin{aligned}
\text{Fermi} &\xrightarrow{\text{mem BW}} \text{Volta} \xrightarrow{\text{GEMM}} \text{Turing} \\
&\xrightarrow{\text{inference}} \text{Ampere} \xrightarrow{\text{sparsity}} \text{Hopper} \\
&\xrightarrow{\text{transformers}} \text{Blackwell} \xrightarrow{\text{scale}} \text{Rubin} \xrightarrow{\text{energy}} \text{?}
\end{aligned}
\end{equation}

Each generation's solution creates the next generation's bottleneck.

\subsection{Precision Trajectory}

\begin{equation}
\text{FP32} \to \text{FP16} \to \text{INT8} \to \text{TF32/BF16} \to \text{FP8} \to \text{FP4 (NVFP4)}
\end{equation}

Precision halves every $\sim$3 years, justified by neural networks' noise tolerance, quantization-aware training, and per-layer adaptive precision. NVFP4 in Rubin represents the current frontier, with further reduction requiring careful validation.

\subsection{Specialization and Memory Bandwidth Scaling}

The ratio of general-purpose FP32 throughput to specialized Tensor Core throughput has declined monotonically. On Volta, FP32 TFLOPS were $\sim$20\% of FP16 Tensor TFLOPS, while on Hopper, that ratio dropped to $\sim$15\%, indicating a clear trend toward specialization for ML workloads.

\begin{table}[h]
\centering
\footnotesize
\begin{tabular}{lcc}
\toprule
Generation & Bandwidth & Growth \\
\midrule
Fermi (2010) & 177 GB/s & --- \\
Volta (2017) & 900 GB/s & 5.1x \\
Ampere (2020) & 2.039 TB/s & 2.2x \\
Hopper (2022) & 3.35 TB/s & 1.7x \\
Blackwell (2024) & 8 TB/s & 2.4x \\
Rubin (2026) & 22 TB/s & 2.75x \\
\bottomrule
\end{tabular}
\caption{Memory bandwidth evolution ($\sim$2x every 2-3 years).}
\end{table}

\subsection{Architectural Invariants}

Despite radical specialization, core structures persist: SM-based organization (16 to 144 SMs), 32-thread warps (unchanged since G80), three-level memory (registers $\to$ shared/L1 $\to$ L2 $\to$ DRAM), and the thread/block/grid model (extended with clusters in Hopper). These invariants provide programming model stability across all seven generations.

\section{Conclusions and Future Directions}

This survey examined GPU architecture evolution through a bottleneck-driven lens, revealing clear causal relationships between workload pressure and architectural innovation.

\subsection{Key Findings}

\textbf{1. Bottleneck-driven evolution}: Each generation directly addresses the dominant bottleneck of its era---memory bandwidth (Fermi), matrix multiplication (Volta), sparsity (Ampere), transformers (Hopper), inference scale (Blackwell).

\textbf{2. Specialization trajectory}: GPUs evolved from general-purpose (Fermi: 0\% specialized) to domain-specific (Blackwell: 90\% specialized). This trend will likely continue, with future GPUs resembling neural processing units.

\textbf{3. Precision halving}: Numerical precision halves every $\sim$3 years (FP32 $\to$ FP16 $\to$ FP8 $\to$ FP4), justified by ML workloads' noise tolerance. NVFP4 (Rubin, 2026) represents the current frontier at 50 PFLOPS inference performance.

\textbf{4. Memory-compute co-scaling}: Memory bandwidth grew $\sim$2x every 2--3 years, tracking compute growth. Future bottlenecks will arise if memory scaling slows.

\textbf{5. Architectural invariants}: Despite specialization, core structures (SM organization, SIMT execution, memory hierarchy) persist, providing programming continuity.

\subsection{The Next Bottleneck: Energy Efficiency}

Our analysis predicts the next fundamental bottleneck: \textbf{energy efficiency}. As data centers approach megawatt-scale power consumption and trillion-parameter models require exaflop training, performance-per-watt becomes the limiting factor. Future architectures may require:
\begin{itemize}
\item Novel device physics (photonics, superconductors)
\item Algorithmic co-design (sparse MoE, state space models)
\item System-level optimization (liquid cooling, renewable energy)
\end{itemize}

\subsection{Implications for Future Research}

\textbf{Hardware-algorithm co-design}: Blackwell's structured sparsity and Hopper's Transformer Engine demonstrate co-design's power. Future algorithms should target architectural primitives (e.g., designing attention variants exploiting hardware features).

\textbf{Precision diversity}: Rather than uniform precision, future models may use adaptive precision per-layer, per-token, or per-parameter---requiring architectural support for fine-grained precision control.

\textbf{Beyond transformers}: If state space models (Mamba, S4) or other architectures replace transformers, hardware must adapt. Architectural agility becomes valuable.

\section*{Acknowledgments}

This survey synthesizes information from NVIDIA architecture whitepapers, technical blogs, and academic publications on deep learning systems.

\bibliography{references}
\bibliographystyle{icml2026}

\end{document}
